\begin{theorem}[Uniform recognition above the threshold]\label{thm:recognition-gap}
There is a number \(\varepsilon_*>0\), independent of smooth type,
such that for every closed smooth homotopy four-sphere \(M\),
\[
               I^+(M)\le T+\varepsilon_*
               \quad\Longrightarrow\quad M\cong_{\rm diff} S^4 .
\]
In particular the same inclusive test applies to the energy of any
one actual positive Yamabe class.
\end{theorem}

\begin{proof}
If nonstandard spheres admit positive classes arbitrarily close to
\(T\), choose such classes on \(M_i\), allowing the type to vary,
with \(e_i\to T\) and physical volume-one Yamabe metrics \(g_i\).
Put \(L=\max\{\lambda_{\max}(W^+),\lambda_{\max}(W^-)\}\) and
\(B=4e/3-\int L^2dV\). Trace freeness gives the pointwise bound
\(L^2\le\tfrac23(|W^+|^2+|W^-|^2)\), so \(B\ge0\).
The strict and equality classifications of Chen and Zhu give
\(Y^2<36\int L^2dV\) on any nonstandard smooth homotopy sphere.
Simple connectivity excludes the nontrivial connected-sum factors,
and the exceptional projective plane and compact simply connected
product of surfaces have nonzero second homology
\cite[Theorems 1.4 and 1.6]{ChenZhu2014Conformal}.
Consequently \eqref{eq:budgets} gives
\[
 D_i+3B_i<8(e_i-T),\qquad Y_i\to y_*=8\sqrt3\,\pi,\qquad
 6\|L_i\|_2\to y_* .
\]

In each class take the volume-one modified Yamabe metric \(q_i\)
with \(R_{q_i}-6L_{q_i}=c_i\le0\)
\cite[Lemma 2.1 and Theorem 1.7]{ChenZhu2014Conformal}.
The constant test gives
\(Y_i\le c_i+6\int L_{q_i}dV_{q_i}\le c_i+6\|L_i\|_2\).
Thus \(c_i\to0\), the variance of \(L_{q_i}\) tends to zero, and
\(R_{q_i}\to y_*\) in \(L^2(q_i)\). The conformally invariant
nonnegative spectral deficit has mass \(B_i\to0\), so the sum of
the two squared Weyl densities tends to \(2T\) in \(L^1(q_i)\).
Gauss--Bonnet and the curvature decomposition now give
\[
 \|\Ric_{q_i}-(y_*/4)q_i\|_2\to0,\qquad
 \big\||\Rm_{q_i}|_{\rm tens}^2-64\pi^2\big\|_1\to0 .
\]
The positive conformal factors admit smooth \(C^2\) approximations
at each index. Take volume-one approximations preserving these limits
with \(R\ge-1\), using \(R_{q_i}\ge c_i\). The ordinary Yamabe
inequality absorbs the \(L^2\) scalar error and supplies a uniform
inhomogeneous Sobolev inequality.

For clarity, the compactness step uses Ricci flow of these smooth
auxiliary metrics. Write \(U=|\Rm|_{\rm tens}\), let
\(\phi'=a\phi^2\), \(\phi(0)=8\pi\), with the dimensional
curvature evolution constant \(a\), and set \(v=(U-\phi)_+\).
Initially \(\|v\|_2\to0\). Entropy monotonicity and the
logarithmic Sobolev inequality give on existing short flows a
uniform Sobolev estimate with potential \(R/4\)
\cite[\S3]{Perelman2002}\cite[Theorem 5.5, arXiv version 4]{YeSobolevFlow}.
Since \(R_+\le C(\phi+v)\), small \(\|v\|_2\) absorbs the potential
into an ordinary Sobolev estimate. Testing
\((\partial_t-\Delta)v\le a(v^2+2\phi v)\) against \(2v\)
then absorbs the cubic term and keeps \(\sup_t\|v\|_2=o(1)\)
on a common short horizon. At a space-time maximum of \(tU\) with
\(U=K\) and \(tK\ge2\), the preceding interval of length \(1/K\)
has curvature at most \(2K\). The curvature gradient estimate
gives \(U\ge K/2\) on a ball of radius \(c/\sqrt K\) at that time.
A cutoff iteration in the ordinary Sobolev inequality bounds its
volume below by \(c'K^{-2}\) for large \(K\). This contradicts
the small tail. The resulting \(U\le C/t\) ensures existence on
a common interval \cite[Theorems 7.1 and 8.1]{HamiltonSingularities}.

Let \(\lambda_0=y_*/4\), \(\lambda'=2\lambda^2\),
\(\lambda(0)=\lambda_0\). In an evolving orthonormal frame
\(P=\Ric-\lambda g\) satisfies
\((\partial_t-\Delta)P=2\Rm(P)+2\lambda P\).
The same Sobolev absorption controls \(\int v|P|^2\), hence
\(\sup_t\|P\|_2=o(1)\). The volume equation
\(V'+4\lambda V=-\int\operatorname{tr}P\,dV\) gives
\(V(t)\to(1-2\lambda_0t)^2\). At fixed small positive time,
bounded curvature and derivatives, ball noncollapse and a diameter
bound from volume packing give a smooth closed limit diffeomorphic
to every sufficiently late source on a subsequence
\cite[Theorem 4.7]{CheegerGromovTaylor1982}\cite{HamiltonCompactness}.
It is positive Einstein with scalar square integral
\(16\lambda(t)^2(1-2\lambda_0t)^2=y_*^2\). By Gauss--Bonnet
and signature its chiral energy is exactly \(T\).
Proposition~\ref{prop:recognition} identifies the standard sphere,
so the round metric pulls back to a zero-energy positive class on each
late source. This contradicts their smooth types. The descent paper
gives the complete smoothing estimates
\cite[section on Ricci flow at the recognition endpoint]{Baker2026FixedDescent}.

This excludes nonstandard positive classes with energy below
\(T+\eta\) for some universal \(\eta>0\). Taking
\(0<\varepsilon_*<\eta\) also covers an infimum equal to
\(T+\varepsilon_*\), by its defining approximation.
\end{proof}
